\chapter*{Acknowledgements}
\setlength{\parskip}{1em}

Agradeço ao meu orientador, Professor António José Borba Ramires Fernandes, pelo acompanhamento atento nesta última etapa académica e pela paciência, conhecimento, bom humor e exemplo profissional que me deu desde a licenciatura. Levo comigo a sua exigência e a sua paixão pela computação gráfica.

Às pessoas mais importantes da minha vida, obrigado. Ao meu pai, Rui, que sempre me apoiou e me mostrou o que é a integridade e força. À minha irmã, Carolina, a voz da razão que me acolheu nos momentos mais difíceis. E à minha mãe, Paula, o meu maior amor, que me trouxe calma e sempre foi a minha fã número um. Só cheguei aqui porque lutámos juntos, e esta conquista é tanto vossa como minha. Agradeço também aos familiares que acompanharam esta viagem, mas, se começasse a escrever nomes, nunca mais daqui saía.

Aos meus amigos, companheiros de batalha passados e futuros, obrigado por tudo o que me ensinaram. Ao Carlos António, o irmão que nunca tive desde o 2.º ano da licenciatura: que venham muitos mais anos. À Márcia Teixeira e à Sofia Marques, as minhas amigas lindas, que prezo mais do que ninguém: nunca saberei retribuir o que fizeram por mim. Agradecer à Bea, ao Miguel, à Filipa, à Morgana, ao Artur, ao Dinis, ao Rui e ao Zé. Sois todos uma bênção que levo no coração para sempre.

Dedico esta dissertação a cinco pessoas. Quatro já não estão entre nós, mas sinto as suas mãos nos meus ombros enquanto escrevo. À minha avó Lia, que continua ao meu lado e me ensinou a ser casmurro e a persistir. Ao meu avô Nelinho, que me inspirou a deixar sempre um bom nome, porque o mundo precisa de mais pessoas boas. Ao meu avô Pimenta, que me ensinou a ser ambicioso e a apontar para as estrelas, por mais impossível que pareça. À minha avó Dina, um dos grandes amores da minha vida, que me deu amor incondicional e cuja falta não se descreve por palavras. E à minha querida Zinda, a minha “criança grande”, a estrelinha dos meus olhos, que me ensinou a viver com leveza, alegria e pela qual sinto falta todos os dias. Amo-vos aos cinco e daria tudo para estarmos todos juntos mais uma vez.

\setlength{\parskip}{0em}
